% Figure: the two members of §3, the Gaussian and the Matern family, at unit canonical scale.
% Referenced from §3.6 after prop:two-members. Generated by scripts/make-fig-members.py from the
% transforms alone (numerical Fourier inversion of (1 + omega^2)^{-gamma}; no Bessel function),
% so the picture rests on prop:two-members' transform and on nothing cited. The tail statements
% in the caption are what the computed curves show, not theorems of this paper.
\begin{figure}[t]
\centering
\includegraphics[width=\textwidth]{fig-members.png}
\caption{The two members of Proposition~\ref{prop:two-members} at unit canonical scale,
computed from their transforms. Left, linear scale: the variance-gamma densities with
transform $(1 + \omega^2)^{-\gamma}$ for $\gamma = \tfrac12, 1, 2, 4$, and a Gaussian of
variance $4$, the variance the $\gamma = 2$ member has by
Proposition~\ref{prop:matern-exponent}(3); the variances of the four densities differ,
$2\gamma$ each. At $\gamma = \tfrac12$ the density is unbounded at the origin and the curve is
clipped; at $\gamma = 1$ it is the Laplace density $e^{-|x|}/2$, with a cusp. Right, the same
curves on a logarithmic scale: the four densities fall off exponentially, the Gaussian as
$e^{-x^2/8}$.}
\label{fig:members}
\end{figure}
